%% ma: L12-B01-0020
%% lop: 12
%% chuong: 1
%% bai: 1
%% dang: 12.1.8
%% loai: TN4
%% muc: TH
%% dap_an: "A"
%% nguon: "(NBV)-ĐỀ SỐ 11-MỖI NGÀY 1 ĐỀ THI 2025.pdf (D11, MỖI NGÀY 1 ĐỀ - Đề số 11), Phần I Câu 6"
%% kien_thuc: "Bài 1 (tính đơn điệu của hàm số bậc ba)"
%% trang_thai: chinh-thuc
%% ghi_chu: "Đáp án gốc A đúng. Cùng ý tưởng với L12-B01-0013 (xét tính đơn điệu hàm bậc ba qua đạo hàm, khác hàm số và cách hỏi), nhập vào dạng đã duyệt."
\begin{ex}
Hàm số $y=-x^3+3x^2-2$ nghịch biến trên khoảng
\choice{\True $(-\infty;0)$}{$(0;2)$}{$(0;+\infty)$}{$(-\infty;2)$}
\loigiai{$y'=-3x^2+6x=-3x(x-2)$; $y'=0\Leftrightarrow x=0$ hoặc $x=2$. Do $a=-1<0$ nên $y'<0$ trên $(-\infty;0)$ và $(2;+\infty)$, $y'>0$ trên $(0;2)$. Hàm số nghịch biến trên $(-\infty;0)$ và $(2;+\infty)$. Chọn A.}
\end{ex}
